\section{Methodology: Expert Interviews}

Wie in den theorischen Grundlagen beschreiben, hat die Interaktion innerhlab von Meetings viele sehr
untershciedlcihe Factten und ebneen (psychologische Effekte, persänlichkeitsmerkmale, Kommunikation)
Kommunkiation im allgeminen zwichen verschiendne Menschen. Über Remotemmeting und mit einem KI-Agent etnsteht
noch eine zusätzliche ebene der Computer Human interaktion. Um diese Kompexität einzufangen sollen eine
qunatitiaven Studie gezielt für diese Arbeit relevanten informationen erheben. Es bestand die hoffung
die komplexität diese zusammenspiels auf wesentliche Merkmale und Erkenntnis zu reduzieren.
Hauuptziel der interviws war es somit den Übergang von der allgmeien recherche bezüglihc meeting über die
hinzunhame der konkreten praixs erfarhung hinzu relevaten Features für einen KI-Agenten zu schaffen.

Semi-structured expert interviews were chosen as the method. This type of interview allows for in-depth
insights into the respondents’ perspectives through open-ended questions and flexible follow-ups, while at the same time
ensuring a thematic focus through a guide.

Expert interviews are particularly suitable because experts, through their specialized knowledge and
practical experience, provide efficient access to specific areas of knowledge \cite[p.
877]{baurHandbuchMethodenEmpirischen2022}. The definition of an expert is context-dependent and must be
determined on a case-by-case basis depending on the research topic \cite[p.
877]{baurHandbuchMethodenEmpirischen2022}. In this study, only individuals with verifiable professional
experience in the field of facilitation and who are professionally active in this context were interviewed.
Therefore, the interviews conducted can be classified as semi-structured expert interviews.

Letzlich ist zu erwarten, dass das der Protpy stellenweise daruaf abziel eine facilitator zu emulieren.
Facilitatoren wenden das "What / When / Who" Framewokr aus [Quelle], welches zum Desing von KI-Systemem von X empfohlen
wird per Defintion ihrer Jobbeschrebigung bereits an. Eine beobachtung aus dieser Praxis stellt somit
eine erweiterte Perspetkive zur Literatur dar.

Es wurde sich auch dazu entschieden Fragen dazu zustellen, wei sich facilitatoren ein solches KI-gestüztes
system vorstellen würden. Man könnte kritisieren, dass Facilitaotren hierfür keine experten sind und es somit
zu spekulativen uninformierten antowrten kommt, da ihnen die technischen möglichkeiten und limitierungen
ki-gestützer systme nciht bewusst sind.
Allerdings sind die fragen so gestellt das der Facilitaotr die verhaltensweisen und
eigenschaften diese systems beschreiben kann, wozu dieser qualifiziert ist. ohne über detailwissen zu der
technsischen implementation verfügen zu müssen. (vgl. Annahnd Frage x.X)
Zustätzlich wurde in Houtti et al. mit Meeting teilnehmden ein co-host design, wie sich meeeting teinnehmedne
ein socheles system wünschen. Hier wird nun eine andere Prespektive geboten. Ein user-cetnric desing
interivew hätte frü diese pespektive vermutlich einen bilden Fleck, weil standard meeting teinmehnden nicht
wissen welche methodkien und facilitating ansaätz zum einsatz kommen und somit auch kein System diebezüglich
designen werden, welches eben jene methodiken und ansaätze berücksichtigt. (Wissen von dem sie nciht wissen, dass
sie es nicht wissen)

Sowohl die Reihenfolge der unterschiedlichen Abschnitte als auch der Fragen innerhalb der
Abschnitte orientieren sich an einer Struktu die vom allgmeinen zum spezieifschen vorgeht.
Es handelt sich um offengestellte Fragen.
Desweitern sind die fragen so gestellt um verhlatensbezogene anker zu triggern um den befrageten zurugen wie
er sich in vergangen situation verhalten hat. [Quelle auch "The mum test"] (Aussagen sind verlössiger wenn
man frage wie wurde das isher gelöst, anstatt wie lässt sich das lösen)

Es hat ein Probe-interview mit einem anderen Studneten stattgefunden, um die allgemeine verständlichekti der frage zu
prüfen und den groben zeitrahne zu prüfen. Da der Stunde allerdings fachfemd war wurde, ein weiter druchlauf
mit CHatGPT druchgeführt. Dabei warude chat gpt als facilitating experte gepromtent und wusste, dass nun ein
Experteninterview simulliert wird. Auch diese interview hatte lediglich das ziel
den Fragebogen zu testen und zielte nicht auf einen erkentnnissgewinn abseites dessen ab. [Quelle das es
best-practice ist einen probelauf zu machen]

% ---------------------------------------------------------------------------------------------------------------------
\subsection{Sampling: Purposeful selection of 5 facilitating experts}
Für diese Studie wurden 5 Experten befragt.

% ---------------------------------------------------------------------------------------------------------------------
\subsection{Data Collection: Conduction of semi-structured interviews (30 min)}
Wann wurden die interviews geführt? Inwlechne zeitraum?
Es wurden nie mehr als 2 Interviews pro Tag gefürht. Es gab einen Interviewer, nämlich den Author dieser Arbeit.
Abwohl die Interviews leicht variert haben, haben diese im schnitt ca. 30min gedauert.
Die interviews haben über Zoom stattgefunden udn wurden, um diese später besser zu transiciribere aufgezeichnet.
Jedes interview startet mit einer kurzen Einleitung in der erklärut wrude, was der zweck des interivews it,
wofür die Datenerhoben werden und wie diese verarbeitet werden. Anschließened wurde die Intervieweten
gefragt, ob diese mit dieser Verarbeitung und insbesondere der Aufnahme einverstnaden sind.
Dem haben alle zugestimmt.

Zu beginn bereits kontext zu liefern, hat dabei geholfen klar einzuordnen worum es in dem Interview gehen soll.

Through
these measures and also through the choice of the interview channel, a certain closeness
could be created, which is helpful in obtaining the views and opinions of the interviewee
31
(Meyen et al., 2019).

Nach dem das interview beednet wurde, wurden sich direkt notzien gemacht und erste gedanken, festzuhalten
währende die eindrück nach den interviews noch frisch waren.

% ---------------------------------------------------------------------------------------------------------------------
\subsection{Analysis: Thematic coding}

Deductive Coding cann miss out on valuable insights. While inductive coding is better for investigated new
ideas or conepts.
Da allerdings bereits eine Literature recherche vorrausgesetllt wurde, besteht zumindest ein verständnis
darüber, welche Themenkomplexe in den interviews zu ewarten sind bzw. welche explizit abgedenkt werden sollen
um genauere informationen von den experten zu erlagnen.
Aus diesem grund wurde ein Kategoriserungssystem basierend auf den literaturrecherchen erstellt.
Und daruf hin wurde im anschlus allerdings eine inductives conding der interview ergebnisse druchg eführt.
dieser ansatz sollte es ermöglichen einenklaren Rehmen für die Interviews zu schaffen aber dennoch offengenug
zu sein um neue, noch nicht in der Rechereche identizifierte informationen aufnhemen zu können.

Somit wurde ein hybrider conding ansatz gewühlt. bei dem die codes iterativs verreinert wurden. Auch die
überkategoreien wurden dann im nach gang weiterfereinert.
